\section*{\centering ABSTRACT}\addcontentsline{toc}{section}{ABSTRACT}
\hspace{1.27cm}
This thesis presents the design, development, and experimental validation of a vision-based Delta robot pick-and-place system using ROS 2, aimed at providing an affordable and functional robotic solution for educational and industrial applications in developing countries, with a particular focus on Cambodia. The Delta robot, known for its parallel kinematic structure and high-speed pick-and-place capability, is implemented as a complete autonomous system integrating computer vision, real-time motion planning, and CAN bus actuation.

\par
\hspace{1.27cm}
The robot is constructed using three RobStride RS-00 quasi-direct-drive actuators controlled via CAN bus at 1\,Mbit/s, an aluminum and 3D-printed structural frame with geometric parameters $f=157$\,mm, $e=35$\,mm, $r_f=200$\,mm, and $r_e=400$\,mm, and a pneumatic solenoid gripper. Open-source software frameworks including Robot Operating System 2 (ROS\,2) and Python are employed for motion planning, kinematics computation, and real-time control. The vision system uses an Intel RealSense D455 stereo depth camera providing aligned RGB and depth streams at up to 30\,fps, combined with a HSV colour-space object detection (with YOLO11n 
interchangeable backbone) for real-time object localisation. Detected pixel coordinates are back-projected to 3D camera-frame coordinates using the pinhole model and transformed to the robot base frame via a calibrated rigid-body transform.

\par
\hspace{1.27cm}
The developed robot features three translational degrees of freedom and operates under a medium motion preset with end-effector velocity $v_{\max}=2000$\,mm/s and acceleration $a_{\max}=5000$\,mm/s$^2$. The kinematic model, including forward and inverse kinematics, was derived analytically following the Trossen Robotics formulation and validated through physical testing, achieving FK position errors consistently below 0.5\,mm. A visual end-effector feedforward correction loop tracks the EE laser marker in the camera frame and iteratively corrects the approach position prior to descent, improving pick accuracy across the non-linear workspace. Conveyor belt prediction is implemented using a triangular motion profile to compensate for robot travel time when picking moving objects at a measured belt speed of 12.85\,mm/s.

\par
\hspace{1.27cm}
Overall, this research demonstrates a complete vision-to-motion pipeline for autonomous pick-and-place, contributes a replicable low-cost Delta robot platform, and promotes local capacity building in robotics and automation engineering in Cambodia.
